\lysh{15654}{2}
\begin{alist}
\item Παρατηρούμε ότι $P(2) = 2^3 - 7 \cdot 2 + 6 = 8 - 14 + 6 = 0$, οπότε το 2 είναι ρίζα του $P(x)$ και το $x-2$ είναι παράγοντας του $P(x)$.
\item Κάνουμε τη διαίρεση $P(x) : (x-2)$,
\begin{center}
\begin{tabular}{r|l}
$x^3 + 0x^2 - 7x + 6$ & $x - 2$ \\ \cline{2-2}
$(+) \quad -x^3 + 2x^2 \phantom{{} - 7x + 6}$ & $x^2 + 2x - 3$ \\ \cline{1-1}
$2x^2 - 7x + 6$ & \\
$(+) \quad -2x^2 + 4x \phantom{{} + 6}$ & \\ \cline{1-1}
$-3x + 6$ & \\
$(+) \quad 3x - 6$ & \\ \cline{1-1}
$0$ &
\end{tabular}
\end{center}
Άρα $P(x) = (x-2)(x^2 + 2x - 3)$. Έχουμε ισοδύναμα:
\begin{align*}
P(x) = 0 &\Leftrightarrow \\
(x-2)(x^2 + 2x - 3) = 0 &\Leftrightarrow \\
\begin{cases} x - 2 = 0 \\ \text{ή} \\ x^2 + 2x - 3 = 0 \end{cases} &\Leftrightarrow
\begin{cases} x = 2 \\ \text{ή} \\ x = 1, x = -3 \end{cases}.
\end{align*}
Άρα οι λύσεις της εξίσωσης είναι $x = 1, x = 2, x = -3$.
\end{alist}

